\documentclass[times, a4paper, 12pt, onecolumn, oneside]{article} 

\usepackage[a4paper,left=2.0cm,right=2.0cm,bottom=1.0cm,top=2.0cm,includefoot,includehead]{geometry}
\usepackage{amsmath}
\usepackage{amsthm}
\usepackage{amsfonts}
\usepackage{amssymb}
\usepackage{graphicx} 
\usepackage{fancyhdr}
\usepackage{natbib}
\usepackage{enumitem}
\setlist{nosep}

\pagestyle{fancyplain}
\date{}
\lhead{{\em Immonen (LUT)}}
\chead{Exercises Week 4}
\rhead{\today}

\bibliographystyle{apalike}

\title{ Exercises Week 4}
\author{Veikka Immonen}
\date{\today}

\begin{document}

\maketitle

\section{Continuation of GP regression from previous exercises}

\subsection{Generate the data from the underlying true process}

Generate noisy data to be modeled with
\[
y = Hx + \epsilon
\]
\begin{itemize}
\item $\theta := (\nu, \ell, \sigma) = (2.5, 0.1, 1.0)$
\item $x \sim \mathcal{GP}(0, c_\theta(x, x'))$
\item $c_\theta(x, x') = \sigma^2 \frac{2^{1 - \nu}}{\Gamma(\nu)}\left(\sqrt{2\nu}\frac{||x - x'||}{\ell}\right)^\nu K_\nu\left(\sqrt{2\nu}\frac{||x - x'||}{\ell}\right)$
\item $H$: select operator
\item $\epsilon \sim \mathcal{N}(0, \sigma_\epsilon^2), ~ \sigma_\epsilon = 0.2$ 
\item $n_\text{data} = 51$
\item $n_\text{state} = 1001$
\end{itemize}
Presented in Figure~\ref{fig:data}

\subsection{Sample $x$ and $\ell$}

The full hierarchical model is
\[
p(x, \theta | y) \propto \underbrace{p(y | x, \theta)}_\text{likelihood} \underbrace{p(x | \theta)}_\text{prior} \underbrace{p(\theta)}_\text{hyperprior}
\]
Now we sample $x$ and $\ell$, hence:
\[
p(x, \ell | y) \propto p(y | x) p(x | \ell) p(\ell),
\]
Like previously same Metropolis-within-Gibbs (MWG) -like schema is used again (with pCN in state sampling!).

\subsection{SPDE Approximation}

For the approximation, let $\alpha = \nu + \frac{d}{2}, ~ d = \dim(t)$, and set distance scale $\kappa = \frac{\sqrt{2\nu}}{\ell}$ in the Matérn covariance. Now, the covariance is a solution to SPDE
\[
\begin{aligned}
\tau(\kappa^2 - \Delta)^\frac{\alpha}{2}\mathcal{X} &= \mathcal{W}, & \mathcal{W} &\sim \mathcal{N}(0, I),
\end{aligned}
\]
where
\[
\begin{aligned}
\tau^2 &= \frac{\Gamma(\nu)}{\Gamma(\alpha)(4\pi)^\frac{d}{4}\kappa^{2\nu}\sigma^2}, &
\Delta = \sum_{i=1}^d \frac{\partial^2}{\partial t_i^2}
\end{aligned}
\]
Approximation on a regular $d$-dimensional grid with spacing constants $h=\{h_1, ..., h_d\}$ using finite difference yields
\[
\begin{aligned}
\tau(\kappa^2 - L^{(h)})^\frac{\alpha}{2}x &= w, & w &\sim \mathcal{N}(0, \underbrace{\prod_{i=1}^d \frac{1}{h_i}}_{=\delta} I),
\end{aligned}
\]
where $L$ is a linear operator of second order central difference. Let $A^{(h)}:= \kappa I + L^{(h)}$. Now we get
\[
\begin{aligned}
x &\sim \mathcal{N}(0, Q^{-1}) & Q &= \delta\tau^2 A^\alpha
\end{aligned}
\]
While A is sparse and hence computationally easy to deal with, its inverse would be still dense. Also, $\alpha$ must be a positive integer. Therefore an alternative iterative sampling algorithm is used to prevent the computation of $Q^{-1}$:
\begin{enumerate}
\item $\xi \sim \mathcal{N}(0, I)$
\item if $\alpha$ is odd:
\begin{itemize}
    \item solve $L^T\tilde{x}_0 = \xi, ~ LL^T = A$
    \item $n_\text{iters} = \frac{\alpha - 1}{2}$
\end{itemize}
\item otherwise:
\begin{itemize}
    \item solve $\tilde{x}_0 = \xi$
    \item $n_\text{iters} = \frac{\alpha}{2}$
\end{itemize}
\item for $i = 1,...,n_\text{iters}$:
\begin{itemize}
    \item solve $A\tilde{x}_{i} = \tilde{x}_{i-1}$
\end{itemize}
\item $x = \frac{\tilde{x}_{n_\text{iters}}} {\tau \sqrt{\delta}}$
\end{enumerate}

Moreover, the prior is also computed without constructing full $Q$. The full prior is
\[
p(x | \theta) \propto \exp\left(-\frac{1}{2}\left( x^TQx - \log |Q| + \frac{d}{2}\log2\pi \right) \right),
\]
where
\begin{itemize}
\item $x^T Qx = \delta\tau^2 x^T\underbrace{A...A}_{\times \alpha} x$
\item $\log |Q| = \log |\tau^2 \delta A^\alpha| = 2n\log\tau + n\log\delta + \alpha\log|A|$
\end{itemize}

\subsection{Arrangements}

\begin{itemize}
\item $(\nu_0, \ell_0, \sigma_0, \sigma_{\epsilon0}) = (2.5, 0.1, 1.0, 0.2)$
\item $\beta_x = \beta_\ell = 0.1$
\item log-Gaussian hyperprior: $\log\ell \sim \mathcal{N}(\log 0.2, 1.0^2)$
\item $\text{Chain length} = 10000$
\end{itemize}

\subsection{Sample with full covariance}

With given environment, sampling took approximately 60 seconds.
Results presented in Figure~\ref{fig:1_1_1}.

\subsection{Sample with SPDE covariance}

With given environment, sampling took less than 10 seconds.
Results presented in Figure~\ref{fig:1_1_2}.



\subsection{Remarks}

Utilizing sparse matrices via SPDE approximation in Matérn field sampling provides remarkably efficient approach to sample states and compute probability density values.

Boundary effects were considered by adding new grid points with original spacing beyond original boundaries. If padding is not added, Laplacian operator add first order difference to boudary points, which
equals to adding homogeneous Neumann boundary condition $\nabla x(t) = 0, ~ t \in \partial\Omega$ to the SPDE. padding reduced the effect of the condition in the sampled area.


\section*{Figures}

\begin{figure}[ht]
    \centering
    \includegraphics[width=0.8\textwidth]{./data_and_true.pdf}
    \caption{True process realization and the noisy data.}
    \label{fig:data}
\end{figure}

\begin{figure}[ht]
    \centering
    \includegraphics[width=0.8\textwidth]{./1_1_1_state.pdf}
    \includegraphics[width=0.8\textwidth]{./1_1_1_pchain.pdf}
    \includegraphics[width=0.8\textwidth]{./1_1_1_phist.pdf}
    \caption{Task 1: Results of sampling $x$ (full covariance) and $\ell$}
    \label{fig:1_1_1}
\end{figure}

\begin{figure}[ht]
    \centering
    \includegraphics[width=0.8\textwidth]{./1_1_2_state.pdf}
    \includegraphics[width=0.8\textwidth]{./1_1_2_pchain.pdf}
    \includegraphics[width=0.8\textwidth]{./1_1_2_phist.pdf}
    \caption{Task 1: Results of sampling $x$ (SPDE covariance) and $\ell$}
    \label{fig:1_1_2}
\end{figure}


% \bibliography{../template/references}

\end{document}
